\documentclass[runningheads]{llncs}
\usepackage[T1]{fontenc}
\usepackage{amsmath,amssymb}
\usepackage{booktabs,multirow,array}
\usepackage{graphicx}
\usepackage[hidelinks]{hyperref}
\usepackage{xspace}
\newcommand{\churnAfterECommerce}{16.6}
\newcommand{\churnAfterIranian}{15.6}
\newcommand{\churnAfterTelco}{26.4}
\newcommand{\churnBeforeECommerce}{16.8}
\newcommand{\churnBeforeIranian}{15.7}
\newcommand{\churnBeforeTelco}{26.5}
\newcommand{\citlICLCellTwoCell}{+1.8}
\newcommand{\citlICLCellTwoCellAbs}{1.8}
\newcommand{\citlICLIranian}{-0.1}
\newcommand{\citlICLIranianAbs}{0.1}
\newcommand{\citlICLMaven}{+0.8}
\newcommand{\citlICLMavenAbs}{0.8}
\newcommand{\citlICLOrange}{+2.1}
\newcommand{\citlICLOrangeAbs}{2.1}
\newcommand{\citlICLTelco}{+0.7}
\newcommand{\citlICLTelcoAbs}{0.7}
\newcommand{\citlPFNThreeCellTwoCell}{-0.2}
\newcommand{\citlPFNThreeCellTwoCellAbs}{0.2}
\newcommand{\citlPFNThreeIranian}{-0.1}
\newcommand{\citlPFNThreeIranianAbs}{0.1}
\newcommand{\citlPFNThreeMaven}{+0.2}
\newcommand{\citlPFNThreeMavenAbs}{0.2}
\newcommand{\citlPFNThreeOrange}{0.0}
\newcommand{\citlPFNThreeTelco}{0.0}
\newcommand{\citlPFNTwoCellTwoCell}{+2.0}
\newcommand{\citlPFNTwoCellTwoCellAbs}{2.0}
\newcommand{\citlPFNTwoIranian}{0.0}
\newcommand{\citlPFNTwoMaven}{+0.7}
\newcommand{\citlPFNTwoMavenAbs}{0.7}
\newcommand{\citlPFNTwoOrange}{+0.2}
\newcommand{\citlPFNTwoOrangeAbs}{0.2}
\newcommand{\citlPFNTwoTelco}{+0.6}
\newcommand{\citlPFNTwoTelcoAbs}{0.6}
\newcommand{\conflictGroupsECommerce}{0}
\newcommand{\conflictGroupsIranian}{14}
\newcommand{\conflictGroupsTelco}{18}
\newcommand{\conflictRowsECommerce}{0}
\newcommand{\conflictRowsIranian}{64}
\newcommand{\conflictRowsTelco}{42}
\newcommand{\ctxCellTwoCell}{30,628}
\newcommand{\ctxCitlOrange}{-35.3}
\newcommand{\ctxCitlOrangeAbs}{35.3}
\newcommand{\ctxOrange}{30,000}
\newcommand{\ctxOrangePFNBal}{+4.4}
\newcommand{\ctxOrangePFNBalAbs}{4.4}
\newcommand{\ctxOrangePFNFull}{+14.3}
\newcommand{\ctxOrangePFNFullAbs}{14.3}
\newcommand{\ctxOrangePFNHalf}{+8.2}
\newcommand{\ctxOrangePFNHalfAbs}{8.2}
\newcommand{\ctxOrangePFNpBal}{0.541}
\newcommand{\ctxOrangePFNpFull}{0.011}
\newcommand{\ctxOrangePFNpHalf}{0.201}
\newcommand{\ctxTelco}{4,201}
\newcommand{\dupCellTwoCell}{0}
\newcommand{\dupECommerce}{556}
\newcommand{\dupIranian}{314}
\newcommand{\dupMaven}{0}
\newcommand{\dupOrange}{0}
\newcommand{\dupPctCellTwoCell}{0.0}
\newcommand{\dupPctECommerce}{9.9}
\newcommand{\dupPctIranian}{10.0}
\newcommand{\dupPctMaven}{0.0}
\newcommand{\dupPctOrange}{0.0}
\newcommand{\dupPctTelco}{0.6}
\newcommand{\dupTelco}{40}
\newcommand{\empAbsICLCellTwoCell}{+0.042}
\newcommand{\empAbsICLCellTwoCellAbs}{0.042}
\newcommand{\empAbsICLIranian}{+0.584}
\newcommand{\empAbsICLIranianAbs}{0.584}
\newcommand{\empAbsICLMaven}{+0.032}
\newcommand{\empAbsICLMavenAbs}{0.032}
\newcommand{\empAbsICLOrange}{-0.026}
\newcommand{\empAbsICLOrangeAbs}{0.026}
\newcommand{\empAbsICLTelco}{+0.057}
\newcommand{\empAbsICLTelcoAbs}{0.057}
\newcommand{\empAbsPFNThreeCellTwoCell}{+0.175}
\newcommand{\empAbsPFNThreeCellTwoCellAbs}{0.175}
\newcommand{\empAbsPFNThreeIranian}{+0.597}
\newcommand{\empAbsPFNThreeIranianAbs}{0.597}
\newcommand{\empAbsPFNThreeMaven}{+0.125}
\newcommand{\empAbsPFNThreeMavenAbs}{0.125}
\newcommand{\empAbsPFNThreeOrange}{+0.082}
\newcommand{\empAbsPFNThreeOrangeAbs}{0.082}
\newcommand{\empAbsPFNThreeTelco}{+0.095}
\newcommand{\empAbsPFNThreeTelcoAbs}{0.095}
\newcommand{\empAbsPFNTwoCellTwoCell}{+0.032}
\newcommand{\empAbsPFNTwoCellTwoCellAbs}{0.032}
\newcommand{\empAbsPFNTwoIranian}{+0.543}
\newcommand{\empAbsPFNTwoIranianAbs}{0.543}
\newcommand{\empAbsPFNTwoMaven}{+0.079}
\newcommand{\empAbsPFNTwoMavenAbs}{0.079}
\newcommand{\empAbsPFNTwoOrange}{-0.017}
\newcommand{\empAbsPFNTwoOrangeAbs}{0.017}
\newcommand{\empAbsPFNTwoTelco}{+0.051}
\newcommand{\empAbsPFNTwoTelcoAbs}{0.051}
\newcommand{\empHiICLCellTwoCell}{+0.072}
\newcommand{\empHiICLCellTwoCellAbs}{0.072}
\newcommand{\empHiICLIranian}{+0.690}
\newcommand{\empHiICLIranianAbs}{0.690}
\newcommand{\empHiICLMaven}{+0.120}
\newcommand{\empHiICLMavenAbs}{0.120}
\newcommand{\empHiICLOrange}{+0.013}
\newcommand{\empHiICLOrangeAbs}{0.013}
\newcommand{\empHiICLTelco}{+0.123}
\newcommand{\empHiICLTelcoAbs}{0.123}
\newcommand{\empHiPFNThreeCellTwoCell}{+0.202}
\newcommand{\empHiPFNThreeCellTwoCellAbs}{0.202}
\newcommand{\empHiPFNThreeIranian}{+0.703}
\newcommand{\empHiPFNThreeIranianAbs}{0.703}
\newcommand{\empHiPFNThreeMaven}{+0.215}
\newcommand{\empHiPFNThreeMavenAbs}{0.215}
\newcommand{\empHiPFNThreeOrange}{+0.114}
\newcommand{\empHiPFNThreeOrangeAbs}{0.114}
\newcommand{\empHiPFNThreeTelco}{+0.159}
\newcommand{\empHiPFNThreeTelcoAbs}{0.159}
\newcommand{\empHiPFNTwoCellTwoCell}{+0.061}
\newcommand{\empHiPFNTwoCellTwoCellAbs}{0.061}
\newcommand{\empHiPFNTwoIranian}{+0.650}
\newcommand{\empHiPFNTwoIranianAbs}{0.650}
\newcommand{\empHiPFNTwoMaven}{+0.154}
\newcommand{\empHiPFNTwoMavenAbs}{0.154}
\newcommand{\empHiPFNTwoOrange}{+0.012}
\newcommand{\empHiPFNTwoOrangeAbs}{0.012}
\newcommand{\empHiPFNTwoTelco}{+0.116}
\newcommand{\empHiPFNTwoTelcoAbs}{0.116}
\newcommand{\empICLCellTwoCell}{+0.5}
\newcommand{\empICLCellTwoCellAbs}{0.5}
\newcommand{\empICLIranian}{+7.4}
\newcommand{\empICLIranianAbs}{7.4}
\newcommand{\empICLMaven}{+0.2}
\newcommand{\empICLMavenAbs}{0.2}
\newcommand{\empICLOrange}{-4.5}
\newcommand{\empICLOrangeAbs}{4.5}
\newcommand{\empICLTelco}{+0.6}
\newcommand{\empICLTelcoAbs}{0.6}
\newcommand{\empLoICLCellTwoCell}{+0.009}
\newcommand{\empLoICLCellTwoCellAbs}{0.009}
\newcommand{\empLoICLIranian}{+0.418}
\newcommand{\empLoICLIranianAbs}{0.418}
\newcommand{\empLoICLMaven}{-0.051}
\newcommand{\empLoICLMavenAbs}{0.051}
\newcommand{\empLoICLOrange}{-0.066}
\newcommand{\empLoICLOrangeAbs}{0.066}
\newcommand{\empLoICLTelco}{0.000}
\newcommand{\empLoPFNThreeCellTwoCell}{+0.145}
\newcommand{\empLoPFNThreeCellTwoCellAbs}{0.145}
\newcommand{\empLoPFNThreeIranian}{+0.439}
\newcommand{\empLoPFNThreeIranianAbs}{0.439}
\newcommand{\empLoPFNThreeMaven}{+0.044}
\newcommand{\empLoPFNThreeMavenAbs}{0.044}
\newcommand{\empLoPFNThreeOrange}{+0.054}
\newcommand{\empLoPFNThreeOrangeAbs}{0.054}
\newcommand{\empLoPFNThreeTelco}{+0.015}
\newcommand{\empLoPFNThreeTelcoAbs}{0.015}
\newcommand{\empLoPFNTwoCellTwoCell}{+0.005}
\newcommand{\empLoPFNTwoCellTwoCellAbs}{0.005}
\newcommand{\empLoPFNTwoIranian}{+0.382}
\newcommand{\empLoPFNTwoIranianAbs}{0.382}
\newcommand{\empLoPFNTwoMaven}{-0.004}
\newcommand{\empLoPFNTwoMavenAbs}{0.004}
\newcommand{\empLoPFNTwoOrange}{-0.044}
\newcommand{\empLoPFNTwoOrangeAbs}{0.044}
\newcommand{\empLoPFNTwoTelco}{-0.012}
\newcommand{\empLoPFNTwoTelcoAbs}{0.012}
\newcommand{\empPFNThreeCellTwoCell}{+1.9}
\newcommand{\empPFNThreeCellTwoCellAbs}{1.9}
\newcommand{\empPFNThreeIranian}{+7.6}
\newcommand{\empPFNThreeIranianAbs}{7.6}
\newcommand{\empPFNThreeMaven}{+1.0}
\newcommand{\empPFNThreeMavenAbs}{1.0}
\newcommand{\empPFNThreeOrange}{+14.3}
\newcommand{\empPFNThreeOrangeAbs}{14.3}
\newcommand{\empPFNThreeTelco}{+0.9}
\newcommand{\empPFNThreeTelcoAbs}{0.9}
\newcommand{\empPFNTwoCellTwoCell}{+0.4}
\newcommand{\empPFNTwoCellTwoCellAbs}{0.4}
\newcommand{\empPFNTwoIranian}{+6.9}
\newcommand{\empPFNTwoIranianAbs}{6.9}
\newcommand{\empPFNTwoMaven}{+0.6}
\newcommand{\empPFNTwoMavenAbs}{0.6}
\newcommand{\empPFNTwoOrange}{-3.0}
\newcommand{\empPFNTwoOrangeAbs}{3.0}
\newcommand{\empPFNTwoTelco}{+0.5}
\newcommand{\empPFNTwoTelcoAbs}{0.5}
\newcommand{\holmICLCellTwoCell}{1.000}
\newcommand{\holmICLIranian}{0.009}
\newcommand{\holmICLMaven}{1.000}
\newcommand{\holmICLOrange}{1.000}
\newcommand{\holmICLTelco}{1.000}
\newcommand{\holmPFNThreeCellTwoCell}{$<$0.001}
\newcommand{\holmPFNThreeIranian}{0.003}
\newcommand{\holmPFNThreeMaven}{1.000}
\newcommand{\holmPFNThreeOrange}{0.155}
\newcommand{\holmPFNThreeTelco}{1.000}
\newcommand{\holmPFNTwoCellTwoCell}{1.000}
\newcommand{\holmPFNTwoIranian}{0.011}
\newcommand{\holmPFNTwoMaven}{1.000}
\newcommand{\holmPFNTwoOrange}{1.000}
\newcommand{\holmPFNTwoTelco}{1.000}
\newcommand{\hoursPFNCacheMillion}{1.25}
\newcommand{\maxAbsCitlPFNThree}{0.2}
\newcommand{\maxChurnShift}{0.3}
\newcommand{\maxDupShift}{1.4}
\newcommand{\maxRawRegretNonOrange}{3.1}
\newcommand{\maxRegretShiftNonOrange}{0.9}
\newcommand{\minHolm}{$<$0.001}
\newcommand{\minHolmCase}{TabPFN-3.5 on Cell2Cell}
\newcommand{\nAllSeedsICL}{1}
\newcommand{\nAllSeedsPFNFast}{2}
\newcommand{\nAllSeedsPFNThree}{3}
\newcommand{\nAllSeedsPFNTwo}{1}
\newcommand{\nCellTwoCell}{51,047}
\newcommand{\nECommerce}{5,074}
\newcommand{\nHierEMPICL}{2}
\newcommand{\nHierEMPPFNFast}{4}
\newcommand{\nHierEMPPFNThree}{5}
\newcommand{\nHierEMPPFNTwo}{2}
\newcommand{\nHolmSig}{6}
\newcommand{\nIranian}{2,836}
\newcommand{\nMaven}{6,589}
\newcommand{\nNegEMPICL}{0}
\newcommand{\nNegEMPPFNFast}{0}
\newcommand{\nNegEMPPFNThree}{0}
\newcommand{\nNegEMPPFNTwo}{0}
\newcommand{\nOrange}{50,000}
\newcommand{\nPrimary}{20}
\newcommand{\nRuns}{685}
\newcommand{\nSeeds}{10}
\newcommand{\nTelco}{7,003}
\newcommand{\nbpICLCellTwoCell}{0.086}
\newcommand{\nbpICLIranian}{$<$0.001}
\newcommand{\nbpICLMaven}{0.699}
\newcommand{\nbpICLOrange}{0.457}
\newcommand{\nbpICLTelco}{0.215}
\newcommand{\nbpPFNThreeCellTwoCell}{$<$0.001}
\newcommand{\nbpPFNThreeIranian}{$<$0.001}
\newcommand{\nbpPFNThreeMaven}{0.097}
\newcommand{\nbpPFNThreeOrange}{0.011}
\newcommand{\nbpPFNThreeTelco}{0.119}
\newcommand{\nbpPFNTwoCellTwoCell}{0.136}
\newcommand{\nbpPFNTwoIranian}{$<$0.001}
\newcommand{\nbpPFNTwoMaven}{0.192}
\newcommand{\nbpPFNTwoOrange}{0.474}
\newcommand{\nbpPFNTwoTelco}{0.370}
\newcommand{\rateCellTwoCell}{28.8}
\newcommand{\rateECommerce}{16.6}
\newcommand{\rateIranian}{15.6}
\newcommand{\rateMaven}{28.4}
\newcommand{\rateOrange}{7.3}
\newcommand{\rateTelco}{26.4}
\newcommand{\regICLOrangeIso}{7.5}
\newcommand{\regICLOrangePlatt}{6.3}
\newcommand{\regICLOrangeRaw}{15.7}
\newcommand{\replDefaultPR}{0.906}
\newcommand{\replICLDefault}{+6.1}
\newcommand{\replICLDefaultAbs}{6.1}
\newcommand{\replICLEMP}{+4.8}
\newcommand{\replICLEMPAbs}{4.8}
\newcommand{\replICLSeeds}{5}
\newcommand{\replICLTuned}{+4.0}
\newcommand{\replICLTunedAbs}{4.0}
\newcommand{\replICLTunedP}{0.023}
\newcommand{\replPFNThreeDefault}{+7.8}
\newcommand{\replPFNThreeDefaultAbs}{7.8}
\newcommand{\replPFNThreeEMP}{+8.3}
\newcommand{\replPFNThreeEMPAbs}{8.3}
\newcommand{\replPFNThreeSeeds}{5}
\newcommand{\replPFNThreeTuned}{+5.8}
\newcommand{\replPFNThreeTunedAbs}{5.8}
\newcommand{\replPFNThreeTunedP}{0.010}
\newcommand{\replTuned}{CatBoost}
\newcommand{\replTunedPR}{0.927}
\newcommand{\seedsPosICLCellTwoCell}{9}
\newcommand{\seedsPosICLIranian}{10}
\newcommand{\seedsPosICLMaven}{8}
\newcommand{\seedsPosICLOrange}{4}
\newcommand{\seedsPosICLTelco}{9}
\newcommand{\seedsPosPFNThreeCellTwoCell}{10}
\newcommand{\seedsPosPFNThreeIranian}{10}
\newcommand{\seedsPosPFNThreeMaven}{9}
\newcommand{\seedsPosPFNThreeOrange}{10}
\newcommand{\seedsPosPFNThreeTelco}{8}
\newcommand{\seedsPosPFNTwoCellTwoCell}{8}
\newcommand{\seedsPosPFNTwoIranian}{10}
\newcommand{\seedsPosPFNTwoMaven}{9}
\newcommand{\seedsPosPFNTwoOrange}{2}
\newcommand{\seedsPosPFNTwoTelco}{8}
\newcommand{\strongDeltaCellTwoCell}{0.0}
\newcommand{\strongDeltaIranian}{-0.8}
\newcommand{\strongDeltaIranianAbs}{0.8}
\newcommand{\strongDeltaOrange}{+0.1}
\newcommand{\strongDeltaOrangeAbs}{0.1}
\newcommand{\strongPFNThreeCellTwoCell}{+1.9}
\newcommand{\strongPFNThreeCellTwoCellAbs}{1.9}
\newcommand{\strongPFNThreeIranian}{+7.6}
\newcommand{\strongPFNThreeIranianAbs}{7.6}
\newcommand{\strongPFNThreeOrange}{+14.3}
\newcommand{\strongPFNThreeOrangeAbs}{14.3}
\newcommand{\strongRefCellTwoCell}{LightGBM (200 trials)}
\newcommand{\strongRefIranian}{LightGBM}
\newcommand{\strongRefOrange}{LightGBM}
\newcommand{\tauEMPCellTwoCell}{0.88}
\newcommand{\tauEMPIranian}{0.79}
\newcommand{\tauEMPMaven}{0.67}
\newcommand{\tauEMPOrange}{0.83}
\newcommand{\tauEMPTelco}{0.52}
\newcommand{\tauEPthrCellTwoCell}{0.41}
\newcommand{\tauEPthrMaven}{0.04}
\newcommand{\tauEPthrTelco}{0.28}
\newcommand{\timeICLCacheOrange}{17}
\newcommand{\timeLGBMOrange}{0.8}
\newcommand{\timePFNCacheOrange}{45}
\newcommand{\timePFNUncachedSmallOrange}{1,686}
\newcommand{\topEMPCellTwoCell}{0}
\newcommand{\topEMPIranian}{6}
\newcommand{\topEMPMaven}{7}
\newcommand{\topEMPOrange}{1}
\newcommand{\topEMPTelco}{7}

\makeatletter\newcommand{\tabin}[1]{\@@input #1 }\makeatother

\setlength{\emergencystretch}{2.5em}
\newcommand{\tfm}{TFM\xspace}
\newcommand{\tfms}{TFMs\xspace}

\begin{document}

\title{Do Tabular Foundation Models Pay Off for Churn?\\
A Profit-, Calibration- and Tuning-Aware Re-evaluation}
\titlerunning{Do Tabular Foundation Models Pay Off for Churn?}

\author{Ismail Hossain Polas}
\authorrunning{I.~H.~Polas}
\institute{Jahangirnagar University\\ \email{ismailhossainpolash@gmail.com}}

\maketitle

\begin{abstract}
Tabular foundation models (\tfms) predict from a labelled context in a single forward pass. A recent benchmark
reported that they decisively outperform tree ensembles for customer churn, but it compared them with untuned
gradient-boosted trees (GBDTs) and scored only ranking metrics, while a retention campaign is judged in money. We
re-evaluate TabPFN-2.5, TabPFN-3.5, TabPFN-3.5-Fast and TabICL v2 against Optuna-tuned LightGBM, XGBoost and CatBoost
on five public telecom churn datasets. Every model is scored on expected maximum profit (EMP), e-Profits, calibration
and scoring time, over five seeds with paired-bootstrap intervals. Only TabPFN-3.5 earns a significant EMP gain on
more than one dataset (Orange \empPFNThreeOrangeAbs\%, Cell2Cell \empPFNThreeCellTwoCellAbs\%, Iranian
\empPFNThreeIranianAbs\%), and it is never significantly worse. TabICL v2, the earlier benchmark's leader, is
significant only on Iranian and underpredicts Orange churn by \citlICLOrangeAbs{} points on a \rateOrange\% base rate.
Profit metrics reorder models mainly when they are statistically tied. A class-balanced context lowers TabPFN-3.5's
Orange gain rather than raising it. On the earlier benchmark's headline dataset, tuning the baseline cuts TabICL v2's
PR-AUC lead from \replICLDefaultAbs{} to \replICLTunedAbs{} points without closing it. Scoring cost is minor; the
licence of TabPFN-3.5's weights, which forbids production use, is the binding constraint.
\keywords{Customer churn \and Tabular foundation models \and Expected maximum profit \and Calibration \and Benchmarking}
\end{abstract}

\section{Introduction}

Churn models exist to choose whom a retention campaign contacts. A model is useful to an operator to the extent that
the customers it ranks highest are the ones worth an incentive, so the natural yardstick is campaign profit rather than
classification accuracy. Verbraken et al.\ made this precise with the expected maximum profit (EMP) criterion
\cite{verbraken2013emp}, and later work showed that selecting churn models by AUC instead of profit leaves money on the
table \cite{verbeke2012newinsights,hoppner2020proftree,manzoor2025eprofits}.

Tabular foundation models change what a churn model is. TabPFN \cite{hollmann2025tabpfn,priorlabs2025tabpfn25,priorlabs2026tabpfn35}
and TabICL \cite{qu2025tabicl,qu2026tabiclv2} are pretrained on synthetic tables and predict a new customer by
conditioning on the labelled training set as context, with no fitting and no hyperparameter search. On general
benchmarks they now lead or match tuned GBDTs \cite{erickson2025tabarena,qu2026tabiclv2,priorlabs2026tabpfn35}. For
churn specifically, Seyedzadeh and Karimi \cite{seyedzadeh2026churn} benchmarked six \tfms on nine Kaggle datasets.
They concluded that \tfms ``collectively and decisively outperform'' classical, deep and tree-ensemble models, with
TabICL v2 ahead of XGBoost by up to 9.23 PR-AUC points.

That study was the first to bring \tfms to churn, and it leaves three questions open. Its XGBoost and random forest
used library-default hyperparameters \cite[Sect.~3.2]{seyedzadeh2026churn}, so part of the reported gap may be a tuning
gap. Its metrics were F1, ROC-AUC, PR-AUC and log loss; its financial estimate scaled the PR-AUC gap by a retention rate
rather than evaluating a campaign. Finally, its own discussion reads TabICL v2's lead as ``operational \ldots\ rather
than as a statistically decisive metric lead'' \cite[Sect.~4.2]{seyedzadeh2026churn}, which leaves the size of the lead
unsettled.

We re-evaluate the question with the baselines tuned and the decision made explicit. Four \tfms (TabPFN-2.5,
TabPFN-3.5, TabPFN-3.5-Fast, TabICL v2) face Optuna-tuned LightGBM, XGBoost and CatBoost, an explainable boosting
machine and logistic regression on five public telecom datasets, three of them from real operators. Every model is
scored on ranking metrics, EMP, e-Profits, a fixed-threshold campaign policy, calibration and scoring time, over five
seeded splits with paired-bootstrap intervals. Our findings:
\begin{enumerate}
\item \textbf{Only TabPFN-3.5 pays off consistently.} Against the best tuned GBDT its EMP gain is significant on
  \nSigEMPPFNThree{} of 5 datasets and never negative beyond noise; TabICL v2 and TabPFN-2.5 are significant only on
  Iranian. On the two fictional datasets, Telco and Maven, no \tfm differs significantly from the GBDTs
  (Sect.~\ref{sec:rq1}).
\item \textbf{Profit metrics reorder ties, not clear winners.} Where one model leads clearly (Orange, Cell2Cell), the
  EMP winner equals the AUC winner in every seed; where models are within noise, the EMP winner differs in up to
  \topEMPMaven{} of 5 seeds (Sect.~\ref{sec:rq2}).
\item \textbf{Calibration matters for one model on one dataset.} Raw TabICL v2 underpredicts Orange churn by
  \citlICLOrangeAbs{} points and loses \regICLOrangeRaw\% of attainable campaign profit at a fixed threshold; elsewhere
  recalibration changes that loss by at most \maxRegretShiftNonOrange{} points (Sect.~\ref{sec:rq3}).
\item \textbf{A balanced context does not help when the full training set fits.} It lowers TabPFN-3.5's Orange gain
  from \ctxOrangePFNFullAbs\% to \ctxOrangePFNBalAbs\% (Sect.~\ref{sec:rq4}).
\item \textbf{Tuning shrinks but does not erase the headline gap.} On the earlier benchmark's E-Commerce dataset,
  TabICL v2's PR-AUC lead falls from \replICLDefaultAbs{} points over default XGBoost to \replICLTunedAbs{} over tuned
  \replTuned{}, and remains significant (Sect.~\ref{sec:rq5}).
\end{enumerate}
We also report scoring throughput and licensing (Sect.~\ref{sec:rq6}), which together decide whether a gain can be
deployed at all. The code, every per-split prediction and the table generator are released
(Sect.~\ref{sec:availability}).

\section{Related Work}

\paragraph{Profit-aware churn evaluation.}
Verbraken et al.\ defined EMP as the expected best campaign profit when the probability that a contacted churner
accepts the offer is uncertain and Beta distributed \cite{verbraken2013emp}. Verbeke et al.\ used profit to compare
churn classifiers in telecom \cite{verbeke2012newinsights}, and Höppner et al.\ optimised EMP directly when growing
decision trees \cite{hoppner2020proftree}. Manzoor et al.\ proposed e-Profits, which replaces EMP's population-level
customer value with a per-customer value derived from Kaplan--Meier retention, and found that it reranks six classical
models on IBM Telco and Maven \cite{manzoor2025eprofits}. None of these studies includes a foundation model.

\paragraph{Tabular foundation models.}
TabPFN v2 \cite{hollmann2025tabpfn} established prior-fitted transformers for tables up to roughly 10k rows;
TabPFN-2.5 extended this to 50k rows \cite{priorlabs2025tabpfn25}, and TabPFN-3 and 3.5 to a million rows with a cached
context for fast scoring \cite{priorlabs2026tabpfn3,priorlabs2026tabpfn35}. TabICL separates column and row attention to
scale in-context learning \cite{qu2025tabicl}; TabICL v2 reports beating a tuned and ensembled RealTabPFN-2.5 on TabArena
and TALENT without tuning \cite{qu2026tabiclv2}. On TabArena, GBDTs remain strong on larger datasets while \tfms excel on
small ones \cite{erickson2025tabarena}, and TabPrep shows that feature engineering, absent from most benchmarks,
changes these comparisons \cite{tschalzev2026tabprep}. TabPFN's predictive probabilities are reported to be well
calibrated out of the box \cite{tanna2025tabtune}, while ensembling \tfms can damage calibration
\cite{tanna2026ensembling}.

\paragraph{\tfms in business prediction.}
The closest analogue to our study is the credit-risk review of Baesens et al.\ \cite{baesens2026creditrisk}. Across
14 probability-of-default datasets TabPFN-3 ranks first, but its margin over the strongest GBDT is narrow and ``driven
by the latest TabPFN and TabICL generations rather than by the model class as a whole''. TabPFN-3 also underpredicts
the mean default rate. Tanna et al.\ found that, under heavy imbalance, building a class-balanced context adds 3 to 4 AUC
points over uniform sampling for subsampled contexts of 1k to 50k rows \cite{tanna2026resampling}; V4FinBench reports
a similar benefit of imbalance-aware context construction for bankruptcy prediction \cite{kostrzewa2026v4finbench}.
For churn, the only \tfm benchmark is \cite{seyedzadeh2026churn}. Table~\ref{tab:gap} sets the closest studies against
ours.

\begin{table}[t]
\caption{Closest prior studies. Tuned: baselines hyperparameter-tuned. Calibration: analysed beyond reporting log
loss. Tests: uncertainty or significance reported for differences between models.}
\label{tab:gap}
\centering\small
\resizebox{\textwidth}{!}{%
\begin{tabular}{@{}llccccc@{}}
\toprule
Study & Domain & \tfms & Tuned & Profit metric & Calibration & Tests \\
\midrule
Manzoor et al.\ \cite{manzoor2025eprofits} & churn & no & yes & e-Profits, EMP & no & yes \\
Baesens et al.\ \cite{baesens2026creditrisk} & credit risk & yes & yes & no & yes & yes \\
Seyedzadeh \& Karimi \cite{seyedzadeh2026churn} & churn & yes & no & no & log loss only & no \\
\midrule
This work & churn & yes & yes & EMP, e-Profits, policy & yes & yes \\
\bottomrule
\end{tabular}}
\end{table}

\section{Benchmark Design}

\subsection{Datasets}
Table~\ref{tab:data} lists the six datasets. Five form the main benchmark. Cell2Cell, Iranian and Orange come from
real operators: Cell2Cell from a US wireless carrier \cite{cell2cell}, Iranian from an Iranian operator through the UCI
repository \cite{iranian}, and Orange from the KDD Cup 2009 churn task \cite{guyon2009kddcup}, which we load from
OpenML (dataset 1112). IBM Telco \cite{telco} and Maven \cite{maven} are fictional teaching datasets; we keep them
because prior churn studies use them, including both \cite{manzoor2025eprofits} and \cite{seyedzadeh2026churn}, and
because a gap between real and fictional data would itself be informative. E-Commerce \cite{ecommerce}, the dataset on
which \cite{seyedzadeh2026churn} reported its largest gain, serves only as a replication check (Sect.~\ref{sec:rq5}) and
is excluded from all cross-dataset summaries.

\begin{table}[t]
\caption{Datasets after preprocessing. Dup.: exact duplicate feature rows removed. e-Prof.: tenure and monthly revenue
available for e-Profits.}
\label{tab:data}
\centering\small
\resizebox{\textwidth}{!}{%
\begin{tabular}{@{}llrrrrcc@{}}
\toprule
Dataset & Source & Rows & Feat. & Churn \% & Dup. & e-Prof. & Role \\
\midrule
\tabin{tables/datasets.tex}
\bottomrule
\end{tabular}}
\end{table}

All models share one preprocessing pipeline. Categorical columns are ordinal-encoded with missing values left as
missing, and columns that are more than 99\% missing or constant are dropped. For Maven we drop customers whose status
is ``Joined'', whose outcome is not yet observed, and the two columns that record the churn reason after the fact. We
remove exact duplicate feature rows, keeping the first: an in-context learner can otherwise look a test customer up in
its own context. This removed \dupTelco{} rows (\dupPctTelco\%) from Telco, \dupIranian{} (\dupPctIranian\%) from
Iranian and \dupECommerce{} (\dupPctECommerce\%) from E-Commerce; the other datasets had none.

\subsection{Models}
\textbf{Foundation models.} TabPFN-2.5, TabPFN-3.5 and TabPFN-3.5-Fast run through the \texttt{tabpfn} package (v9.2.0) with default settings and the training-size guardrail disabled; TabICL v2 runs through \texttt{tabicl} (v2.2.0, checkpoint of 2026-02-12). Each
receives the full training split as its context, at the natural churn rate. None is tuned, which is how these models
are meant to be used.

\textbf{Tuned baselines.} LightGBM \cite{ke2017lightgbm} and XGBoost \cite{chen2016xgboost} are tuned with 40 trials
of Optuna's TPE sampler \cite{akiba2019optuna}, and CatBoost \cite{prokhorenkova2018catboost} with 25, each under a
20-minute cap. Trials are scored by validation AUC on an inner 80/20 split of the training data with early stopping;
the best configuration is then refit on the whole training split with the selected number of trees. We also report an
explainable boosting machine (EBM) \cite{nori2019interpretml} and logistic regression with one-hot encoding and median
imputation, both at defaults.

\subsection{Protocol}
For each dataset and each of five seeds we draw a stratified 60/20/20 split into training, calibration and test sets.
All models see identical splits, which the analysis verifies by comparing test indices across runs. Models are fitted
(or given context) on the training split. The calibration split is used only to fit post-hoc recalibration, and every
reported metric is computed on the test split. Uncertainty for a difference between two models is a paired percentile
bootstrap: we resample test customers 500 times per seed, recompute the metric for both models on each resample, and
pool the 2,500 differences across seeds. We call a difference significant when its 95\% interval excludes zero. Each
\tfm is compared with the best tuned GBDT for that dataset and metric, chosen by mean over seeds, which favours the
baseline. GBDTs ran on Kaggle CPU sessions and \tfms on a single NVIDIA T4 (16\,GB).

\subsection{Metrics}
\textbf{Ranking.} ROC-AUC, PR-AUC and lift in the top decile.

\textbf{EMP.} For a campaign that contacts the fraction of customers with the highest scores, a contacted churner
accepts the offer with probability $\gamma$ and is then worth $\mathrm{CLV}-\delta-\phi$; a churner who declines costs
the contact cost $\phi$; a contacted non-churner takes the incentive and costs $\delta+\phi$. EMP is the expected
profit per customer at the profit-maximising cutoff, averaged over $\gamma\sim\mathrm{Beta}(6,14)$
\cite{verbraken2013emp}. We use the standard parameters $\mathrm{CLV}=200$, $\delta=10$ and $\phi=1$, and integrate over
$\gamma$ on a 200-point grid. Because EMP chooses the best cutoff for each model, it depends only on the ranking and
cannot reward calibration.

\textbf{e-Profits.} Following \cite{manzoor2025eprofits}, each customer's value is
$\mathrm{CLV}_i = R_i M/(1-\min(r_i,0.995))$, with monthly revenue $R_i$, margin $M=0.3$, and one-month retention
$r_i = S(T_i+1)/S(T_i)$ from a Kaplan--Meier curve $S$ \cite{kaplan1958nonparametric} fitted on the training split at
the customer's tenure $T_i$. The offer costs $0.1\,\mathrm{CLV}_i$ and the contact $\max(5, 0.3\cdot\text{offer})$. A
contacted churner yields $\mathrm{CLV}_i$ minus both costs; a contacted non-churner yields minus both costs. We report
e-Profits when contacting customers with predicted probability at least 0.5, as in \cite{manzoor2025eprofits}, and when
contacting the top 20\%. It is defined only where tenure and revenue exist (Table~\ref{tab:data}).

\textbf{Campaign policy and calibration.} An operator who acts on probabilities contacts customer $i$ when its expected
profit under the EMP cost model is positive at the mean acceptance rate $\bar\gamma=0.3$. Unlike EMP, this policy uses
the probability itself, so miscalibration costs money. We report its regret: the profit lost relative to the best
cutoff on the same ranking, as a share of that best profit, as the median over a $3\times3$ grid of
$\mathrm{CLV}\in\{100,200,500\}$ and $\delta\in\{5,10,20\}$. Calibration is measured by expected calibration error
(ECE, 15 equal-width bins), Brier score, and calibration-in-the-large: the observed churn rate minus the mean predicted
probability, in percentage points, where a positive value means churn is underpredicted. Each model is evaluated raw and
after two recalibrations fitted on the calibration split: isotonic regression \cite{zadrozny2002transforming} and Platt
scaling \cite{platt1999probabilistic}.

\section{Results}

\subsection{Do foundation models beat tuned GBDTs?}\label{sec:rq1}

TabPFN-3.5 is the only \tfm with a significant EMP gain on more than one dataset (Table~\ref{tab:boot}). Against the
best tuned GBDT it gains \empPFNThreeOrangeAbs\% on Orange, \empPFNThreeCellTwoCellAbs\% on Cell2Cell and
\empPFNThreeIranianAbs\% on Iranian, and it is never significantly worse. TabICL v2 and TabPFN-2.5 are significant only
on Iranian, where every boosted and foundation model scores above 0.98 AUC (Appendix, Table~\ref{tab:auc}) and the
remaining errors are few.
TabPFN-3.5-Fast sits between its parent and TabPFN-2.5.

The Orange result carries the largest effect and the weakest certainty. The \empPFNThreeOrangeAbs\% gain corresponds to
an EMP difference whose 95\% interval starts at \empLoPFNThreeOrange{} per customer, just above zero, because Orange
has only \rateOrange\% churners. We therefore read it as evidence of a real gain of uncertain size.

On the two fictional datasets the picture is flat. No \tfm differs significantly from the best GBDT on Telco or Maven,
on any metric; TabPFN-3.5's EMP gains there are \empPFNThreeTelcoAbs\% and \empPFNThreeMavenAbs\%. The version also matters
more than the family: TabPFN-2.5 and TabPFN-3.5 share an architecture lineage, yet only the latter separates from the
GBDTs on Orange and Cell2Cell. This matches the pattern Baesens et al.\ report for credit risk
\cite{baesens2026creditrisk}.

\begin{table}[t]
\caption{Mean PR-AUC and EMP (profit per customer, CLV units) over five seeds; best per column in bold.}
\label{tab:main}
\centering\scriptsize\setlength{\tabcolsep}{3.2pt}
\resizebox{\textwidth}{!}{%
\begin{tabular}{@{}lcccccccccc@{}}
\toprule
& \multicolumn{2}{c}{Telco} & \multicolumn{2}{c}{Maven} & \multicolumn{2}{c}{Cell2Cell} & \multicolumn{2}{c}{Iranian} & \multicolumn{2}{c}{Orange} \\
\cmidrule(lr){2-3}\cmidrule(lr){4-5}\cmidrule(lr){6-7}\cmidrule(lr){8-9}\cmidrule(l){10-11}
Model & PR-AUC & EMP & PR-AUC & EMP & PR-AUC & EMP & PR-AUC & EMP & PR-AUC & EMP \\
\midrule
\tabin{tables/main_results.tex}
\bottomrule
\end{tabular}}
\end{table}

\begin{table}[t]
\caption{Each \tfm minus the best tuned GBDT: PR-AUC difference in points ($\times100$) and EMP difference relative to
the GBDT. Bold with $\ast$: the 95\% paired-bootstrap interval excludes zero. The reference is the best tuned GBDT
per dataset and metric, shown as PR-AUC / EMP reference (LGBM: LightGBM, XGB: XGBoost, CatB: CatBoost).}
\label{tab:boot}
\centering\scriptsize\setlength{\tabcolsep}{3pt}
\resizebox{\textwidth}{!}{%
\begin{tabular}{@{}lcccccccccc@{}}
\toprule
& \multicolumn{2}{c}{Telco} & \multicolumn{2}{c}{Maven} & \multicolumn{2}{c}{Cell2Cell} & \multicolumn{2}{c}{Iranian} & \multicolumn{2}{c}{Orange} \\
\cmidrule(lr){2-3}\cmidrule(lr){4-5}\cmidrule(lr){6-7}\cmidrule(lr){8-9}\cmidrule(l){10-11}
& \tabin{tables/bootstrap_refs.tex}
\tfm & $\Delta$PR & $\Delta$EMP & $\Delta$PR & $\Delta$EMP & $\Delta$PR & $\Delta$EMP & $\Delta$PR & $\Delta$EMP & $\Delta$PR & $\Delta$EMP \\
\midrule
\tabin{tables/bootstrap.tex}
\bottomrule
\end{tabular}}
\end{table}

\subsection{Do profit metrics change which model wins?}\label{sec:rq2}

They change it when models are tied and leave it alone otherwise (Table~\ref{tab:rank}). On Orange and Cell2Cell,
where TabPFN-3.5 leads clearly, the EMP winner is the AUC winner in all five seeds, and Kendall's $\tau$ between the two
rankings is \tauEMPOrange{} and \tauEMPCellTwoCell. Telco and Maven behave differently: their top models are within noise
of each other, and the EMP winner differs from the AUC winner in \topEMPTelco{} and \topEMPMaven{} of 5 seeds.

e-Profits at a 0.5 threshold disagrees with AUC much more ($\tau=\tauEPthrMaven$ on Maven). This is expected from its
construction rather than evidence of a better ranking: at a fixed threshold the metric also rewards a model whose
probabilities happen to sit on the right side of 0.5, so it mixes ranking with calibration. The reranking reported for
e-Profits on Telco and Maven \cite{manzoor2025eprofits} is therefore reproduced here, and our bootstrap intervals
locate it among models whose ranking metrics are statistically indistinguishable.

\begin{table}[t]
\caption{Agreement between the AUC ranking of all nine models and profit-based rankings: mean Kendall $\tau$ over
seeds, and seeds in which the top model changes. e-Profits needs tenure and revenue.}
\label{tab:rank}
\centering\small
\begin{tabular}{@{}lcccccc@{}}
\toprule
& \multicolumn{2}{c}{EMP} & \multicolumn{2}{c}{e-Profits, top 20\%} & \multicolumn{2}{c}{e-Profits, $p\ge0.5$} \\
\cmidrule(lr){2-3}\cmidrule(lr){4-5}\cmidrule(l){6-7}
Dataset & $\tau$ & top-1 changed & $\tau$ & top-1 changed & $\tau$ & top-1 changed \\
\midrule
\tabin{tables/rank_agreement.tex}
\bottomrule
\end{tabular}
\end{table}

\subsection{Does calibration cost money?}\label{sec:rq3}

Mostly it does not. Outside Orange, the fixed-threshold policy loses at most \maxRawRegretNonOrange\% of attainable
profit for any model without recalibration, and isotonic or Platt scaling changes that loss by at most
\maxRegretShiftNonOrange{} points (Table~\ref{tab:cal}).

Orange is the exception, and one model drives it. Raw TabICL v2 predicts \citlICLOrangeAbs{} percentage points too
little churn on a base rate of \rateOrange\%, and the policy built on those probabilities loses \regICLOrangeRaw\% of
attainable profit. Isotonic regression cuts the loss to \regICLOrangeIso\% and Platt scaling to \regICLOrangePlatt\%.
TabICL v2 also underpredicts on Cell2Cell (by \citlICLCellTwoCellAbs{} points), as does TabPFN-2.5
(\citlPFNTwoCellTwoCellAbs{}). TabPFN-3.5's calibration-in-the-large stays within $\pm$\maxAbsCitlPFNThree{} points on
every dataset, so it is the only \tfm here that can be thresholded without recalibration. The tuned GBDTs are close to
zero as well, consistent with their fitting to the observed base rate.

\begin{table}[t]
\caption{Calibration. CITL: observed minus mean predicted churn rate (points; positive = underprediction). Regret:
share of best-cutoff profit lost by the fixed-threshold policy (median over the cost grid), raw / isotonic / Platt.
The first row per dataset is the best tuned GBDT.}
\label{tab:cal}
\centering\scriptsize
\begin{tabular}{@{}llrrrrrr@{}}
\toprule
& & & \multicolumn{2}{c}{ECE} & \multicolumn{3}{c}{Policy regret (\%)} \\
\cmidrule(lr){4-5}\cmidrule(l){6-8}
Dataset & Model & CITL & raw & iso. & raw & iso. & Platt \\
\midrule
\tabin{tables/calibration.tex}
\bottomrule
\end{tabular}
\end{table}

\subsection{Does a balanced context help?}\label{sec:rq4}

Under heavy imbalance, a balanced context has been reported to add 3 to 4 AUC points for \tfms
\cite{tanna2026resampling}. If that held here, our full-context runs would understate the \tfms on Orange. We kept every
churner in the training split and sampled one or two non-churners per churner, for TabPFN-3.5 and TabICL v2 on Orange
and Cell2Cell, on the same splits and seeds.

Balancing hurt (Table~\ref{tab:ctx}). TabPFN-3.5's EMP gain on Orange fell from \ctxOrangePFNFullAbs\% with the full
context to \ctxOrangePFNHalfAbs\% at 1:2 and \ctxOrangePFNBalAbs\% at 1:1, and neither of the latter is significant.
TabICL v2 gained nothing on either dataset. Balancing also inflated raw predicted churn: on Orange the 1:1 context
overpredicts the churn rate by \ctxCitlOrangeAbs{} points, an error that isotonic recalibration on the calibration
split removes. The difference from \cite{tanna2026resampling} is the regime: there, contexts were subsampled to at most 50k
rows from larger data, so balancing chose which customers to keep. Here the whole training split fits, and
undersampling only discards non-churners.

\begin{table}[t]
\caption{Context construction. Full: the whole training split at its natural churn rate. 1:2 and 1:1: all churners
plus two or one sampled non-churners per churner. $\Delta$EMP is relative to the best tuned GBDT (bold $\ast$:
significant). CITL in points, raw.}
\label{tab:ctx}
\centering\scriptsize
\begin{tabular}{@{}lllrrrrr@{}}
\toprule
Dataset & Model & Context & AUC & PR-AUC & EMP & $\Delta$EMP & CITL \\
\midrule
\tabin{tables/context.tex}
\bottomrule
\end{tabular}
\end{table}

\subsection{Does the earlier headline survive tuning?}\label{sec:rq5}

We reran our protocol on E-Commerce, where \cite{seyedzadeh2026churn} reported TabICL v2 9.23 PR-AUC points ahead of
XGBoost, adding a library-default XGBoost to reproduce their baseline (Table~\ref{tab:repl}). The direction survives.
TabICL v2 leads default XGBoost by \replICLDefaultAbs{} points and the best tuned GBDT, \replTuned, by
\replICLTunedAbs{} points, both significant, and its EMP gain is \replICLEMPAbs\%. Tuning raised the GBDT's PR-AUC
from \replDefaultPR{} to \replTunedPR{} and so removed about a third of the gap. TabPFN-3.5 leads by more
(\replPFNThreeTunedAbs{} points over tuned \replTuned; EMP gain \replPFNThreeEMPAbs\%). Our protocol differs from theirs (a 60/20/20 holdout over five seeds against
five-fold cross-validation), so we compare default and tuned baselines within our own runs rather than with their
numbers. On this dataset, then, \tfms win clearly. The origin of E-Commerce is not documented, so it cannot tell us
whether the effect is specific to real or synthetic data.

\begin{table}[t]
\caption{Replication on E-Commerce. PR-AUC mean (sd) over five seeds; $\Delta$ in PR-AUC points versus default XGBoost
and versus the best tuned GBDT, and EMP relative to the best tuned GBDT (bold $\ast$: significant).}
\label{tab:repl}
\centering\scriptsize
\begin{tabular}{@{}lccccc@{}}
\toprule
Model & PR-AUC & EMP & $\Delta$PR vs default & $\Delta$PR vs tuned & $\Delta$EMP \\
\midrule
\tabin{tables/replication.tex}
\bottomrule
\end{tabular}
\end{table}

\subsection{Scoring cost and licensing}\label{sec:rq6}\label{sec:cost}

An in-context learner re-reads its context for every scoring call unless the context is cached, so its cost depends on
how scoring is run. We timed scoring of 1k, 10k and 50k customers resampled from the test split, after a warm-up call,
on one T4 (Table~\ref{tab:time}). With the context cached, TabPFN-3.5 scores 10k Orange customers in
\timePFNCacheOrange{}\,s against a \ctxOrange-row context, and TabICL v2 with its key-value cache in
\timeICLCacheOrange{}\,s, against \timeLGBMOrange{}\,s for a 1,000-tree LightGBM on CPU. A monthly scoring run of a
million customers therefore costs about \hoursPFNCacheMillion{} T4-hours with cached TabPFN-3.5. That is slow next to
a GBDT but small next to the per-customer profit differences in Table~\ref{tab:boot}. Without caching, small batches
are very expensive (\timePFNUncachedSmallOrange\,s per 10k for TabPFN-3.5 at 1k rows on Orange), and TabPFN-2.5 was
impractical at 30k-row contexts.

Licensing constrains deployment more than compute does. The TabPFN-3.5 weights are released for research and
evaluation, but ``the models, their derivatives, and their outputs cannot be used for commercial or production
purposes'', including ``using model outputs as inputs to internal commercial decision-making''
\cite[Sect.~4]{priorlabs2026tabpfn35}\cite{tabpfn35license}; production use requires a commercial licence. TabICL v2's
weights are BSD-3 licensed. An operator choosing between them therefore faces a trade-off our benchmark makes concrete:
the model that pays off consistently needs a commercial agreement, and the openly licensed one is not reliably better
than a tuned GBDT on these datasets.

\begin{table}[t]
\caption{Seconds to score 10k customers on one T4 (LightGBM on CPU), best of 10k and 50k batches. $\dagger$: only the
1k batch was run because larger ones exceeded our time budget; per-customer time falls with batch size, so this is an
upper bound. Context: training split (\ctxTelco{} / \ctxCellTwoCell{} / \ctxOrange{} rows).}
\label{tab:time}
\centering\small
\begin{tabular}{@{}lrrr@{}}
\toprule
Model & Telco & Cell2Cell & Orange \\
\midrule
\tabin{tables/timing.tex}
\bottomrule
\end{tabular}
\end{table}

\section{Discussion}

The results support a narrower statement than ``foundation models outperform trees for churn''. Gains are real on some
datasets and absent on others, and three factors explain most of the variation we see. The model version matters more
than the model class: TabPFN-3.5 separates from tuned GBDTs where TabPFN-2.5 does not. Baseline tuning removes part of
any measured lead, about a third of it on E-Commerce. And the datasets most used in churn research, Telco and Maven,
are the ones on which no model can be told apart, so studies built on them can neither confirm nor refute a \tfm
advantage.

For a practitioner this suggests an order of operations. Tune the GBDT first; it is the baseline that matters, and
default settings overstate what a \tfm adds. If TabPFN-3.5 still leads on held-out profit and the licence permits it,
use it with a cached context. If only openly licensed models are possible, test TabICL v2 on one's own data and
recalibrate it before thresholding. A balanced context is not a free improvement when the training set fits.

The finding most likely to change with more evidence is the size of the Orange gain. Its interval is wide, and a
stronger GBDT pipeline with feature engineering \cite{tschalzev2026tabprep} or a larger tuning budget
\cite{erickson2025tabarena} could narrow it. The finding least likely to change is the calibration pattern, which is
consistent across seeds and matches the underprediction reported for credit risk \cite{baesens2026creditrisk}.

\section{Limitations}

Our GBDT tuning is lighter than benchmark practice: 25 to 40 TPE trials, against 200 configurations with post-hoc
ensembling in TabArena \cite{erickson2025tabarena}. Because heavier tuning can only help the baselines, our \tfm gains
are, if anything, overstated. We did no feature engineering, which recent work shows can shift model comparisons
\cite{tschalzev2026tabprep}.

The evidence rests on three real-operator datasets, and the largest effect, on Orange, has an interval close to zero.
We cannot claim that TabPFN-3.5's gain generalises to churn in other industries or to operators with millions of
customers, where the context would no longer hold the full training set. The cost model behind EMP and the policy uses
the standard population-level parameters \cite{verbraken2013emp}; e-Profits gives a customer-level alternative on three
datasets only.

We used one splitting protocol (60/20/20 holdout, five seeds) rather than repeated cross-validation, one GPU type, and
the default inference settings of each \tfm. We did not evaluate the API-only TabPFN-3.5-Plus and -Thinking variants,
or the other \tfms covered by \cite{seyedzadeh2026churn}.

\section{Conclusion}

Measured with tuned baselines and in money, tabular foundation models pay off for churn selectively. TabPFN-3.5
improves campaign profit significantly on three of five telecom datasets and never loses beyond noise, but its weights
cannot be used in production without a commercial licence. TabICL v2 and TabPFN-2.5 are rarely distinguishable from a
tuned GBDT, and TabICL v2 needs recalibration before its probabilities can drive a campaign. Profit metrics confirm
clear winners and reshuffle ties. We release the pipeline and every per-split prediction so that these conclusions can
be re-checked as models and licences change.

\section*{Data and Code Availability}\label{sec:availability}
\begin{sloppypar}
Code and results are at \url{https://github.com/ihpolash/tfm-churn-profit-audit} [TK: version DOI]. The repository
contains the prediction and analysis code as scripts and as notebooks, the Kaggle kernel definitions that ran them, the
per-split predictions and split indices of all \nRuns{} runs, the analysis outputs, and the generator and audit for every
table and number in this paper. The datasets are public (Table~\ref{tab:data}) but are not redistributed, because
several are licensed for non-commercial or attributed use only. The labels needed to recompute metrics from the
predictions are regenerated from the original sources by the released code.
\end{sloppypar}

\section*{Declaration on the Use of Generative AI}
The author used an AI assistant (Claude, Anthropic) for implementation and execution of the experiment code, for
literature search and reference checking, for table generation, and for drafting and editing this manuscript. The
research question was chosen by the author; the benchmark design, evaluation protocol and interpretation of the
results were developed in iterative discussion with the assistant, and every design decision and conclusion was
reviewed and approved by the author. All numerical results were produced by the released code. Every table and in-text
number is generated by a released script from the released analysis outputs, which are computed from the released
per-split predictions, and an audit notebook re-derives them before submission. Every reference was checked against its primary source or an authoritative
index. The author takes full responsibility for all content of this paper, irrespective of how it was drafted.

\bibliographystyle{splncs04}
\bibliography{references}

\appendix
\section{Additional Results}

\begin{table}[h]
\caption{Mean ROC-AUC over five seeds; best per column in bold.}
\label{tab:auc}
\centering\small
\begin{tabular}{@{}lccccc@{}}
\toprule
Model & Telco & Maven & Cell2Cell & Iranian & Orange \\
\midrule
\tabin{tables/auc_results.tex}
\bottomrule
\end{tabular}
\end{table}

\end{document}
